\paperchapter{July 2004}{25 July 2004}{C}{50}{GSET 2004}
\keytable{\kcell{1}{D} & \kcell{2}{B} & \kcell{3}{A} & \kcell{4}{D} & \kcell{5}{A} & \kcell{6}{B} & \kcell{7}{D} & \kcell{8}{D} & \kcell{9}{B} & \kcell{10}{C} \\
\kcell{11}{C} & \kcell{12}{C} & \kcell{13}{C} & \kcell{14}{B} & \kcell{15}{D} & \kcell{16}{B} & \kcell{17}{D} & \kcell{18}{D} & \kcell{19}{B} & \kcell{20}{A} \\
\kcell{21}{D} & \kcell{22}{C} & \kcell{23}{C} & \kcell{24}{B} & \kcell{25}{D} & \kcell{26}{C} & \kcell{27}{B} & \kcell{28}{A} & \kcell{29}{B} & \kcell{30}{D} \\
\kcell{31}{A} & \kcell{32}{D} & \kcell{33}{B} & \kcell{34}{D} & \kcell{35}{D} & \kcell{36}{C} & \kcell{37}{B} & \kcell{38}{C} & \kcell{39}{B} & \kcell{40}{C} \\
\kcell{41}{C} & \kcell{42}{D} & \kcell{43}{B} & \kcell{44}{A} & \kcell{45}{B} & \kcell{46}{B} & \kcell{47}{C} & \kcell{48}{A} & \kcell{49}{D} & \kcell{50}{B}}
\begin{context}
\textbf{Directions (Q.~Nos.~1 to 5):} Read the para:\par\smallskip
Some matters do not decompose. Some matters take longer time to decompose. All plants decompose. All creepers are plants. Plastics decompose in longer time. Animals decompose in longer time.\par\smallskip
\textbf{Which of the following statement(s) \emph{strictly} do not flow from above given statements?} (Q.~1--5)
\end{context}
\begin{question}{1}
{Which of the following statement(s) strictly do not flow from the above statements?\par}
\optstwo{Plants are matter}{Everything decomposable is matter}{Animals are matters}{All of the above}
\answer{D}{All of the above}
\begin{explain}
Nowhere does the paragraph say that plants, animals or ``everything decomposable'' are \emph{matter}. It only says \emph{some} matters decompose. So (A), (B) and (C) all fail to follow strictly: \textbf{all of the above}.
\tip{``Strictly follows'' means it must be deducible from the given statements alone, without outside knowledge.}
\end{explain}
\end{question}

\begin{question}{2}
{Which of the following statement(s) strictly do not flow from the above statements?\par}
\optstwo{Creepers decompose}{Creepers are matters}{Some creepers are not animals}{None of the above}
\answer{B}{Creepers are matters}
\begin{explain}
``All creepers are plants'' and ``All plants decompose'' give ``creepers decompose'', which follows. But nothing says plants (or creepers) are \emph{matter}, so \textbf{``creepers are matters'' does not follow}.
\end{explain}
\end{question}

\begin{question}{3}
{Which of the following statement(s) strictly do not flow from the above statements?\par}
\optsone{Plastics are matters}{What takes longer time to decompose is matter}{Creepers decompose}{All of the above}
\answer{A}{Plastics are matters}
\begin{explain}
``Plastics decompose in longer time'' does not tell us plastics are \emph{matter}, so \textbf{``plastics are matters''} does not strictly follow (answer as per the key). ``Creepers decompose'' does follow.
\end{explain}
\end{question}

\begin{question}{4}
{Which of the following statement(s) strictly do not flow from the above statements?\par}
\optsone{What decomposes in longer time is matter}{Matters are creepers}{Animals are matters}{All of the above}
\answer{D}{All of the above}
\begin{explain}
``Some matters take longer time'' does not mean \emph{whatever} takes longer is matter (A); ``matters are creepers'' (B) and ``animals are matters'' (C) are not stated either. \textbf{All of the above} do not follow.
\end{explain}
\end{question}

\begin{question}{5}
{Which of the following statement(s) strictly do not flow from the above statements?\par}
\optstwo{Some matters are plants}{Plastics take time to decompose}{Some matters decompose, others do not}{None of the above}
\answer{A}{Some matters are plants}
\begin{explain}
``Plastics take time to decompose'' (B) is stated. ``Some matters decompose, others do not'' (C) follows from the first two sentences. \textbf{``Some matters are plants''} (A) does not follow, since plants are never linked to matter.
\end{explain}
\end{question}

\begin{context}
\textbf{Directions (Q.~Nos.~6 to 10):} The following is the data (in crores) for educational expenditure (Ed.\ Ex.) for the various plan periods, population (Pop) and Price Index (PI) [Base 1973 $=$ 100]. This is purely hypothetical.
\begin{center}\small
\begin{tabular}{lcccc}\toprule \textbf{Items} & \textbf{VI} & \textbf{VII} & \textbf{VIII} & \textbf{IX}\\\midrule
Ed.\ Ex. & 17.3 & 28.5 & 37.6 & 38.5\\ Pop. & 50 & 55 & 61 & 70\\ PI & 123 & 144 & 189 & 210\\\bottomrule
\end{tabular}
\end{center}
Worked values used below:
\begin{center}\small
\begin{tabular}{lcccc}\toprule & \textbf{VI} & \textbf{VII} & \textbf{VIII} & \textbf{IX}\\\midrule
Increase in Ed.\ Ex.\ over previous plan & -- & 64.7\% & 31.9\% & 2.4\%\\
Per capita Ed.\ Ex.\ (Ed.Ex./Pop) & 0.346 & 0.518 & 0.616 & 0.550\\
Per capita at constant prices ($\times100/$PI) & 0.281 & 0.360 & 0.326 & 0.262\\
Population increase over previous plan & -- & 5 (10\%) & 6 (10.9\%) & 9 (14.8\%)\\\bottomrule
\end{tabular}
\end{center}
\end{context}
\begin{question}{6}
{In which plan period the total increase (\%) in Educational expenditure was maximum?\par}
\optsfour{VI}{VII}{VIII}{IX}
\answer{B}{VII}
\begin{explain}
From VI to VII: $\dfrac{28.5-17.3}{17.3}\times100=64.7\%$, far higher than VIII (31.9\%) or IX (2.4\%). The answer is \textbf{VII}.
\end{explain}
\end{question}

\begin{question}{7}
{In which plan the per capita expenditure (without correction for Pop \& PI) was the maximum?\par}
\optsfour{VI}{VII}{VIII}{IX}
\answer{D}{IX}
\keynote{From the table the maximum per capita expenditure is in Plan VIII (0.616), not IX (0.550). The official key gives (D).}
\begin{explain}
Per capita expenditure $=$ Ed.\ Ex.\ $\div$ Pop: VI 0.346, VII 0.518, \textbf{VIII 0.616}, IX 0.550. By the data, Plan VIII is highest; the key's answer is IX. Mark (D) as per the official key.
\end{explain}
\end{question}

\begin{question}{8}
{In which plan the population increase with respect to the preceding plan period was minimum?\par}
\optsfour{VI}{VII}{VIII}{IX}
\answer{D}{IX}
\keynote{By the data the smallest increase is in Plan VII (5 crore, 10\%). The official key gives (D).}
\begin{explain}
Population increases: VII $+5$, VIII $+6$, IX $+9$. The minimum is VII; the key gives IX. Mark (D) as per the official key.
\end{explain}
\end{question}

\begin{question}{9}
{In which plan the per capita expenditure with PI-100 on education was maximum ?\par}
\optsfour{VI}{VII}{VIII}{IX}
\answer{B}{VII}
\begin{explain}
Per capita expenditure corrected for price (PI $=100$): $\dfrac{\text{Ed.Ex.}/\text{Pop}}{\text{PI}}\times100$: VI 0.281, \textbf{VII 0.360}, VIII 0.326, IX 0.262. The maximum is \textbf{VII}.
\end{explain}
\end{question}

\begin{question}{10}
{In which plan educational expenditure as percent of total allocation was maximum ?\par}
\optsfour{VI}{VII}{VIII}{IX}
\answer{C}{VIII}
\keynote{The table does not give the total plan allocation, so this cannot be computed from the data. The key gives (C) VIII.}
\begin{explain}
Expenditure as a percentage of total allocation needs the total outlay of each plan, which is not supplied. The official key answer is \textbf{VIII}.
\end{explain}
\end{question}

\begin{question}{11}
{Which one of the following would be your reaction to the statement "Research leads to wasteful expenditure"?\par}
\optstwo{I agree with the statement}{I am not sure}{I disagree with the statement}{It is true in most of the cases}
\answer{C}{I disagree with the statement}
\begin{explain}
Research produces new knowledge and solutions, so it is an investment, not waste. The appropriate reaction is \textbf{``I disagree with the statement''}.
\end{explain}
\end{question}

\begin{question}{12}
{Which one of the following is most essential for a researcher ?\par}
\optsone{Genuine interest in Research}{Good understanding of Research Process}{Adequate interest in and aptitude for Research}{Patience and Perseverance}
\answer{C}{Adequate interest in and aptitude for Research}
\begin{explain}
The most essential attribute is \textbf{adequate interest in and aptitude for research}, which covers both motivation and ability.
\end{explain}
\end{question}

\begin{question}{13}
{A mature researcher is one\par}
\optsone{Whose hypothesis and thesis are identical}{Who maintains critical gap between hypothesis and thesis}{Who has an open mind about the results}{Whose hypothesis is very often rejected}
\answer{C}{Who has an open mind about the results}
\begin{explain}
A mature researcher keeps an \textbf{open mind about the results}. The hypothesis is tested, not defended, and the findings are accepted whichever way they go.
\end{explain}
\end{question}

\begin{question}{14}
{Which one of the following is the purpose of research?\par}
\optstwo{to create knowledge}{to construct reality}{to solve problems}{All of the above}
\answer{B}{to construct reality}
\begin{explain}
Per the key, the purpose of research is \textbf{to construct reality}: to build an understanding of reality through systematic inquiry.
\end{explain}
\end{question}

\begin{question}{15}
{Which one of the following tools would yield relatively most valid data ?\par}
\optstwo{Attitude Scale}{Rating Scale}{Opinionnaire}{Observation Schedules}
\answer{D}{Observation Schedules}
\begin{explain}
\textbf{Observation schedules} record actual behaviour directly, so they give the most valid data. Attitude scales, rating scales and opinionnaires depend on self-report or judgement.
\end{explain}
\end{question}

\begin{question}{16}
{A person who harbours no absolutes and appreciates logical alternatives adheres to:\par}
\optstwo{Scientific Realism}{Positive Relativism}{Logical Empiricism}{Logical Positivism}
\answer{B}{Positive Relativism}
\begin{explain}
Accepting no absolute truths while appreciating logical alternatives is \textbf{positive relativism}.
\end{explain}
\end{question}

\begin{question}{17}
{'X' includes Y. Y is included in 'Z'. Therefore :\par}
\optstwo{X is included in Y}{Y is included in X}{Both are same}{All the above are possible}
\answer{D}{All the above are possible}
\begin{explain}
``X includes Y'' ($Y\subseteq X$) and ``Y is included in Z'' ($Y\subseteq Z$) tell us nothing about how X and Z relate. Options (A)--(C) could each be true in some case, so \textbf{all are possible}.
\end{explain}
\end{question}

\begin{question}{18}
{You are working on an urgent assignment at night in your study. The assignment has to be submitted the next morning. Suddenly, it is pitch dark in your house. Tick the option that would be your first response to this situation.\par}
\optsone{Ring up the night duty staff of the electricity office}{Check the main fuse}{Curse the electricity board and their supply}{Check up if there is light in the adjacent houses}
\answer{D}{Check up if there is light in the adjacent houses}
\begin{explain}
The first rational step is to \textbf{check whether the adjacent houses have light}. This tells you whether it is a general power cut or a local fault before you act.
\end{explain}
\end{question}

\begin{question}{19}
{The most effective way for a teacher to make his students punctual would be to\par}
\optsone{lecture on it every second day}{be punctual himself}{use teaching aids to make students understand its importance}{narrate stories on benefits of punctuality}
\answer{B}{be punctual himself}
\begin{explain}
A teacher teaches punctuality best by \textbf{being punctual himself or herself}: example is better than precept.
\end{explain}
\end{question}

\begin{question}{20}
{Why are many people eager to join teaching profession ?\par}
\optsone{Teachers serve the entire society}{More job opportunities are available in teaching profession}{Adequate salary is paid in this profession}{Teaching is considered an easy job}
\answer{A}{Teachers serve the entire society}
\begin{explain}
Per the key, many join teaching because \textbf{teachers serve the entire society}: it is a noble, service-oriented profession.
\end{explain}
\end{question}

\begin{question}{21}
{The most desirable quality in a teacher is\par}
\optstwo{Mastery of the subject}{Mastery of teaching competency}{Clarity of thought and expression}{All of the above}
\answer{D}{All of the above}
\begin{explain}
A good teacher needs subject mastery, teaching competence and clarity of thought and expression: \textbf{all of the above}.
\end{explain}
\end{question}

\begin{question}{22}
{Which one of the following can be the most effective method for building student's character?\par}
\optstwo{Lecturing on moral values}{Teaching about the lives of great men}{Living example of a good teacher}{Praising good character}
\answer{C}{Living example of a good teacher}
\begin{explain}
Character is built most effectively by the \textbf{living example of a good teacher} (values are caught, not taught).
\end{explain}
\end{question}

\begin{question}{23}
{A student comes late to the class daily. You as a teacher would\par}
\optsone{pull up the student with sarcastic remarks}{not allow the student to enter your class}{tell the student to be in time}{allow the student to come in quietly}
\answer{C}{tell the student to be in time}
\begin{explain}
Per the key, the teacher should \textbf{tell the student to be in time}: a firm, respectful correction without humiliation or exclusion.
\end{explain}
\end{question}

\begin{question}{24}
{Which of the following is the most undesirable for a teacher ?\par}
\optstwo{To be authoritative in the class}{To humiliate the students}{To keep the students at a distance}{To treat the students as friends}
\answer{B}{To humiliate the students}
\begin{explain}
\textbf{Humiliating students} is the most undesirable behaviour. It damages self-esteem and the learning relationship.
\end{explain}
\end{question}

\begin{question}{25}
{When answering a question in class if a student gives wrong response, the teacher should:\par}
\optsone{correct the student immediately}{not put any question to the student}{give the correct answer himself}{re-direct the same question to another student}
\answer{D}{re-direct the same question to another student}
\begin{explain}
Re-directing the question to another student keeps the class engaged and lets peers correct the error without embarrassing the first student.
\end{explain}
\end{question}

\begin{context}
\textbf{Direction (Q.~Nos.~26 to 30):} Read the following text and answer questions 26 to 30.\par\smallskip
Television was, towards the end of the war in Vietnam, the main source of information for most Americans, and the images on screen, in black and white, changed war reporting forever. Morley Safer's expose of American soldiers burning a village had a huge impact on viewers at home, and added to the military's distrust of reporters. In Vietnam, correspondents were made honorary majors (World War II reporters had the momentary rank of captain) and could travel anywhere on military transport if space were available. The only condition: that they could not betray troop movements, a request that was respected. The press battles were over misleading briefings which ran contrary to what the reporters could see for themselves in the field.

To this day, some military men glumly insist that television and the press lost the war in Vietnam by demoralizing the home front. The U.S.\ military would never again allow such access to its operations as it did in Vietnam.

The tension between the military and the press will never cease, because both need each other but cannot grant the other what it really wants.
\end{context}
\begin{question}{26}
{During the Vietnam war, the military started doubting the press because of;\par}
\optsone{the availability of TV for covering war}{the video-cassettes on the war}{the war images telecast by American television companies}{the television images being black and white}
\answer{C}{the war images telecast by American television companies}
\begin{explain}
Morley Safer's televised exposé of soldiers burning a village ``added to the military's distrust of reporters'', i.e.\ \textbf{the war images telecast by American television companies}.
\end{explain}
\end{question}

\begin{question}{27}
{What was the major difference between war reporters of World War II and those of the Vietnam war? The Vietnam war reporters:\par}
\optsone{could go to all war fronts}{had a higher rank and more facilities}{respected their promise}{had more disagreements with military authorities}
\answer{B}{had a higher rank and more facilities}
\begin{explain}
Vietnam correspondents were made \textbf{honorary majors} (a higher rank than WWII reporters, who were captains) and could travel on military transport: a higher rank and more facilities.
\end{explain}
\end{question}

\begin{question}{28}
{The press battles were the result of:\par}
\optsone{the difference between the versions of military authority and the war reporters}{the different thinking and reporting of different reporters}{the press reporters' different views on troop movement}{the different press reports of various newspapers}
\answer{A}{the difference between the versions of military authority and the war reporters}
\begin{explain}
``The press battles were over misleading briefings which ran contrary to what the reporters could see for themselves'', i.e.\ the \textbf{difference between the military's version and the reporters' version}.
\end{explain}
\end{question}

\begin{question}{29}
{Some military men still believe that America lost the Vietnam war because:\par}
\optsone{the press reporters betrayed the cause}{the mass media made people cry against the war}{the television images on war were improper}{the army became demoralised}
\answer{B}{the mass media made people cry against the war}
\begin{explain}
Some military men insist that ``television and the press lost the war in Vietnam by \textbf{demoralizing the home front}'', which the key paraphrases as the mass media turning people against the war.
\end{explain}
\end{question}

\begin{question}{30}
{The tension between the military and the press will continue because :\par}
\optsone{both of them want to emerge superior}{the press opposes the military}{the press serves the people and the military is faithful to the state}{the press pursues truth in all situations}
\answer{D}{the press pursues truth in all situations}
\begin{explain}
The passage says both ``need each other but cannot grant the other what it really wants''. The key gives (D): the press pursues truth in all situations, which the military cannot always permit.
\end{explain}
\end{question}

\begin{question}{31}
{Mark the word that does not belong to the group :\par}
\optsfour{Smile}{Anger}{Annoyance}{Irritation}
\answer{A}{Smile}
\begin{explain}
Anger, annoyance and irritation are negative emotions of displeasure. \textbf{Smile} is a positive expression, so it does not belong.
\end{explain}
\end{question}

\begin{question}{32}
{Spot the expression that does not belong to the group:\par}
\optstwo{Gnashing teeth}{Knitting brows}{Creased foreheads}{Smiling eyes}
\answer{D}{Smiling eyes}
\begin{explain}
Gnashing teeth, knitted brows and creased foreheads express anger or worry. \textbf{Smiling eyes} express happiness.
\end{explain}
\end{question}

\begin{question}{33}
{Bird: Feather: Cat:?\par}
\optsfour{Claw}{Fur}{Ears}{Skin}
\answer{B}{Fur}
\begin{explain}
A bird's body is covered with \textbf{feathers}; a cat's body is covered with \textbf{fur}.
\end{explain}
\end{question}

\begin{question}{34}
{Heavy: Lead:: Cool:?\par}
\optsfour{Winter}{Water}{Ice}{Breeze}
\answer{D}{Breeze}
\begin{explain}
Lead is a typical example of something \emph{heavy}; a \textbf{breeze} is a typical example of something \emph{cool} (per the key).
\end{explain}
\end{question}

\begin{question}{35}
{Air: Animates:: Vitamins: ?\par}
\optsfour{Amino Acids}{Mineral}{Food}{Health}
\answer{D}{Health}
\begin{explain}
Air animates (gives life); vitamins give \textbf{health}.
\end{explain}
\end{question}

\begin{question}{36}
{The average age of 32 students is 10 years. If the teacher's age is also included, then their average age increases by one year. What is the teacher's age in years ?\par}
\optsfour{33}{21}{43}{45}
\answer{C}{43}
\begin{explain}
Total age of 32 students $=320$. With the teacher, 33 persons average 11, so the total is $363$. Teacher $=363-320=\mathbf{43}$ years.
\end{explain}
\end{question}

\begin{question}{37}
{The average weight of a ring is 10 gm which contains gold and copper in the ratio of 9.5 : 0.5. What would be the weight of copper in the ring weighing 8 gms if the same ratio is to be maintained ?\par}
\optsfour{0.5 gm}{0.4 gm}{0.3 gm}{3.5 gm}
\answer{B}{0.4 gm}
\begin{explain}
Copper fraction $=\dfrac{0.5}{9.5+0.5}=\dfrac{1}{20}$. In an 8 g ring, copper $=\dfrac{8}{20}=\mathbf{0.4\ g}$.
\end{explain}
\end{question}

\begin{question}{38}
{In a garden, there are peacocks and four-footed animals. All these have 98 feet and 35 heads. How many peacocks are there?\par}
\optsfour{23}{25}{21}{19}
\answer{C}{21}
\begin{explain}
Let $p$ = peacocks and $a$ = animals: $p+a=35$ and $2p+4a=98$, so $2a=98-70=28$, $a=14$ and $p=\mathbf{21}$.
\end{explain}
\end{question}

\begin{question}{39}
{Grass looks green because :\par}
\optstwo{The eye responds to the green colour.}{The grass reflects only green colour.}{Green is one of the primary colours.}{Green is not a warm colour.}
\answer{B}{The grass reflects only green colour.}
\begin{explain}
An object looks green because it \textbf{reflects green light} and absorbs the other colours.
\end{explain}
\end{question}

\begin{question}{40}
{In a coded message FRIEND is written as IOKCQA. Using the same code SPIKER will be written as:\par}
\optsfour{VMLIGT}{VNHHIV}{VMKIHO}{VNFMHP}
\answer{C}{VMKIHO}
\begin{explain}
FRIEND $\rightarrow$ IOKCQA: the shifts are $+3, -3, +2, -2, +3, -3$.
SPIKER: S$+3=$V, P$-3=$M, I$+2=$K, K$-2=$I, E$+3=$H, R$-3=$O, giving \textbf{VMKIHO}.
\end{explain}
\end{question}

\begin{question}{41}
{Identify the statement 'Quality of research is dependent on researcher's knowledge of research methodology' as:\par}
\optsfour{Assumption}{Fact}{Hypothesis}{Opinion}
\answer{C}{Hypothesis}
\begin{explain}
``Quality of research depends on the researcher's knowledge of methodology'' is a testable proposition about a relationship, i.e.\ a \textbf{hypothesis}.
\end{explain}
\end{question}

\begin{question}{42}
{Identify the statement 'Teaching is an art' as:\par}
\optsfour{Assumption}{Fact}{Hypothesis}{Opinion}
\answer{D}{Opinion}
\begin{explain}
``Teaching is an art'' is a value judgement: an \textbf{opinion}.
\end{explain}
\end{question}

\begin{question}{43}
{Identify the statement 'Physical growth of a learner is age related' as:\par}
\optsfour{Assumption}{Fact}{Hypothesis}{Opinion}
\answer{B}{Fact}
\begin{explain}
``Physical growth of a learner is age-related'' is an established \textbf{fact}.
\end{explain}
\end{question}

\begin{question}{44}
{Identify the statement 'Manipulation of independent variable/variables is required in experimentation' as:\par}
\optsfour{Assumption}{Fact}{Hypothesis}{Opinion}
\answer{A}{Assumption}
\begin{explain}
Per the key, ``manipulation of independent variables is required in experimentation'' is treated as an \textbf{assumption}, a premise on which experimental research rests.
\end{explain}
\end{question}

\begin{question}{45}
{Identify the statement 'Protein is essential for reconstruction of cells' as:\par}
\optsfour{Assumption}{Fact}{Hypothesis}{Opinion}
\answer{B}{Fact}
\begin{explain}
``Protein is essential for reconstruction of cells'' is a scientific \textbf{fact}.
\end{explain}
\end{question}

\begin{question}{46}
{1, 3, 5, 7, 11, 13, \ldots{}?\ldots{}\par}
\optsfour{15}{17}{19}{21}
\answer{B}{17}
\begin{explain}
The sequence is 1 followed by prime numbers: 1, 3, 5, 7, 11, 13, \textbf{17}.
\end{explain}
\end{question}

\begin{question}{47}
{3, 15, 35, 77, \ldots{}?\ldots{}\par}
\optsfour{142}{156}{141}{172}
\answer{C}{141}
\keynote{The logical next term is $11\times13=143$, which is not offered. The official key gives (C) 141.}
\begin{explain}
Each term is a product of consecutive odd primes: $1\times3=3$, $3\times5=15$, $5\times7=35$, $7\times11=77$, so next $11\times13=143$. The closest option, and the key's answer, is \textbf{141}.
\end{explain}
\end{question}

\begin{question}{48}
{3, 5, 9, 15, 23, 33, 45, \ldots{}\par}
\optsfour{59}{61}{68}{69}
\answer{A}{59}
\begin{explain}
The differences increase by 2 each time: $+2, +4, +6, +8, +10, +12$, so next $+14$: $45+14=\mathbf{59}$.
\end{explain}
\end{question}

\begin{question}{49}
{4, 2, 5, 1, 6, \ldots{}?\ldots{}\par}
\optsfour{8}{7}{9}{0}
\answer{D}{0}
\begin{explain}
Two alternating series: $4, 5, 6$ (increasing) and $2, 1, \mathbf{0}$ (decreasing).
\end{explain}
\end{question}

\begin{question}{50}
{1221, 2442, 3663, 4884, \ldots{}?\ldots{}\par}
\optsfour{6084}{6006}{6000}{5665}
\answer{B}{6006}
\keynote{$5\times1221=6105$ is not among the options; the official key gives (B) 6006.}
\begin{explain}
The terms are $1221\times1, \times2, \times3, \times4$. The next would be $1221\times5=6105$. The key's intended pattern is palindromic numbers ($a\,2a\,2a\,a$), and it gives \textbf{6006}.
\end{explain}
\end{question}
